\documentclass[preprint,12pt]{elsarticle}

\usepackage{amsmath,amssymb}
\usepackage{graphicx}
\usepackage{booktabs}

\journal{Journal of Computational Science}

\begin{document}

\begin{frontmatter}

\title{Residual Transport Maps for Calibrated Scientific Emulators}
\author[inst1]{Elena V. Petrov\corref{cor1}}
\ead{epetrov@example.edu}
\author[inst1,inst2]{Minho Park}
\ead{mpark@example.edu}
\cortext[cor1]{Corresponding author}
\address[inst1]{Department of Applied Mathematics, Example University, Boston, 02115, United States}
\address[inst2]{Center for Scientific Computing, Example University, Boston, 02115, United States}

\begin{abstract}
We describe a compiling \texttt{elsarticle} manuscript for journal submission. A frozen scientific emulator is calibrated with a monotone residual transport map estimated on a small holdout set. The procedure improves interval coverage on three partial-differential-equation benchmarks without retraining the surrogate. The source uses the standard Elsevier front matter, keywords, and numbered bibliography so it typesets with a conventional TeX Live installation that provides \texttt{elsarticle.cls}.
\end{abstract}

\begin{keyword}
scientific machine learning \sep uncertainty quantification \sep elsarticle \sep emulator calibration
\end{keyword}

\end{frontmatter}

\section{Introduction}

Learned surrogates are now common in computational science, but their residual law is rarely calibrated when the operating envelope changes. Journal workflows still expect a familiar \texttt{elsarticle} skeleton: title, authors, affiliations, abstract, keywords, numbered sections, and a compact bibliography.

Let $f_\theta$ denote a pretrained emulator. We treat calibration as estimating a monotone map on a scalar residual score, then inverting that map to form prediction intervals.

\section{Calibration}

For a labeled holdout set $\{(x_i,y_i)\}_{i=1}^{n}$ define $s_i=\|y_i-f_\theta(x_i)\|_2$. Write $\hat F$ for the empirical distribution function of $\{s_i\}$. At a query $x$ the symmetric interval of level $1-\alpha$ is
\begin{equation}
  f_\theta(x)\pm \hat F^{-1}(1-\alpha).
\end{equation}
Exchangeability of calibration and test scores yields the usual conformal $1/(n+1)$ coverage correction.

\section{Related work}

Elsevier journals in computational science typically expect a short account of prior art before the method. Physics-informed networks and operator learning supply the frozen emulator~\cite{raissi2019}. Conformal prediction supplies the coverage argument~\cite{vovk2005}. We treat residual scores as a one-dimensional interface so that neither ingredient has to be rewritten.

\section{Numerical evidence}

Table~\ref{tab:elsevier-coverage} summarizes holdout coverage before and after calibration. Width increases modestly while coverage approaches the $90\%$ target. A second split with $n=400$ changes coverage by less than $0.03$, which suggests that the empirical map is already stable at the smaller holdout size used in the main table.

\begin{table}[t]
\centering
\caption{Holdout coverage and mean interval width for a $90\%$ target.}
\label{tab:elsevier-coverage}
\begin{tabular}{@{}lcc@{}}
\toprule
Problem & Uncalibrated coverage & Calibrated coverage \\
\midrule
Viscous Burgers equation & 0.73 & 0.90 \\
Darcy flow & 0.66 & 0.88 \\
Reaction--diffusion & 0.79 & 0.91 \\
\bottomrule
\end{tabular}
\end{table}

\section{Conclusion}

The manuscript is intentionally short: it demonstrates a complete \texttt{elsarticle} document that compiles, and it records a practical calibration baseline for scientific emulators.

\begin{thebibliography}{9}

\bibitem{raissi2019}
M.~Raissi, P.~Perdikaris, G.~E. Karniadakis, Physics-informed neural networks, J. Comput. Phys. 378 (2019) 686--707.

\bibitem{vovk2005}
V.~Vovk, A.~Gammerman, G.~Shafer, Algorithmic Learning in a Random World, Springer, 2005.

\end{thebibliography}

\end{document}
